(1) Isotropy is general: $v \in W_u$ gives $\ip{v}{v}u = c(v)^2u = 0$, and $v_1, v_2 \in W_u$ give $2\ip{v_1}{v_2}u = \{c(v_1), c(v_2)\}u = 0$, so $W_u$ is totally isotropic, hence contained in its own orthogonal complement, so $2\dim W_u \le 6$ and $\dim W_u \le 3$. For the vacuum $u_0 = 1 \in \Lambda^0 W$: $c(w')1 = 0$ for $w' \in W'$ and $c(w+w')1 = \sqrt2\,w$, so $W_{u_0} = W'$, of dimension exactly $3$. By Proposition \claimref{correspondence.spin-representation} the spin image acts transitively on $\T\setminus\{0\}$, and the construction is equivariant, $W_{gu} = g(W_u)$ (from $c(gv) = g c(v) g^{-1}$), so $\dim W_u = 3$ for every $u \neq 0$. Joint kernel: for $W_{u_0} = W'$ the joint kernel of the three contractions $\iota_{w'}$ is $\Lambda^0 W = \C\cdot 1$ (if $x$ has a nonzero component on a monomial $w_S$ with $S \neq \emptyset$, pick any $i \in S$: the $w_{S\setminus i}$-component of $\iota_{w_i'}x$ comes from $w_S$ alone, since $\iota_{w_i'}w_{S'} = \pm w_{S'\setminus i}$ or $0$ and $S'\setminus i = S\setminus i$ with $i \in S'$ forces $S' = S$, so it is nonzero); equivariance transports the statement to every $u$.
(2) The image $c\bigl(\Lambda^\bullet U'\bigr)u$ equals $c(\Cl(V_\C))u$: the anticommutation relations reorder any word in the basis vectors of $U' \oplus W_u$ into a combination of normally ordered products $c(\xi')c(\eta)$ with $\xi'$ a monomial in $U'$ and $\eta$ a monomial in $W_u$, so these $8\cdot 8 = 64$ products span $\Cl(V_\C)$; acting on $u$, the right-hand factor $c(\eta)$ gives $0$ unless $\eta = 1$, since $c(W_u)u = 0$; $c(\Cl(V_\C))u$ is a nonzero submodule, hence everything, by irreducibility (Theorem \claimref{correspondence.fock-model-complexified-ascended-algebra}). So the map is surjective from a space of dimension $2^3 = 8$ onto a space of dimension $8$, hence bijective. The parity statement holds because $c(\xi)$ for even (odd) $\xi$ preserves (swaps) the chirality halves and $u$ is even; if $u \in \T'$ instead, the roles swap. Normalising a basis of $W_u$ against the dual basis of $U'$ produces three pairs with $f^2 = f^{*2} = 0$, $ff^* + f^*f = 1$: the ladder pairs of Fact \claimref{correspondence.imported-ladder-pairs-fock-modules-over}, here over $\C$; the script checks all CAR relations exactly. Complements dual to $W_u$ exist: for a basis $w_1, w_2, w_3$ of $W_u$ choose $v_j$ with $\ip{w_i}{v_j} = \delta_{ij}$ (the functionals $\ip{w_i}{\cdot}$ are independent) and set $u_j := v_j - \tfrac12\sum_k \ip{v_j}{v_k}\,w_k$; then $\ip{w_i}{u_j} = \delta_{ij}$ and $\ip{u_j}{u_k} = \ip{v_j}{v_k} - \tfrac12\ip{v_k}{v_j} - \tfrac12\ip{v_j}{v_k} = 0$, so $U' := \operatorname{span}(u_j)$ is a dual isotropic complement. Nothing below depends on the choice.
(3) The first sentence is Theorem \claimref{correspondence.fock-model-complexified-ascended-algebra} read through (2): a vector $v = w + w'$ acts as a creation operator plus an annihilation operator of the $u$-adapted system. Disjointness: if $U = W_{u} = W_{u'}$ with $u \in \T\setminus 0$, $u' \in \T'\setminus 0$, the joint kernel of $c(U)$ would contain the two-dimensional space $\C u \oplus \C u'$, contradicting (1).
